% from ega4-11-fr.tex
\oldpage[IV-3]{150}
\begin{proof}
Par hypothèse, l'intersection des noyaux des $u_\alpha$ est réduite à $0$~;
comme $A$ est artinien, il y a déjà un nombre fini de ces noyaux dont
l'intersection est $0$, donc on peut se borner au cas où la famille
$(A'_\alpha)$ est finie. Soit $k'$ une clôture algébrique de $k$~; on sait
$(\mathbf 0_{\mathrm{III}},10.3.1)$ qu'il existe un homomorphisme local
$A\to B$ faisant de $B$ un $A$-module plat, tel que $B$ soit un anneau local
artinien, entier sur $A$ et que $B\otimes_A k$ soit isomorphe à $k'$.

Par platitude, les noyaux des homomorphismes
$B\to B'_\alpha=B\otimes_A A'_\alpha$ se déduisent de ceux des $u_\alpha$ par
tensorisation avec $B$, et comme ils sont en nombre fini, leur intersection est
réduite à $0$ $(\mathbf 0_{\mathrm I},6.1.3)$. Considérons les anneaux
$B'_{\alpha\beta}$, localisés de $B'_\alpha$ en ses idéaux maximaux~; on sait
(Bourbaki, \emph{Alg. comm.}, chap.~II, \S~3, n$^\circ$~3, cor.~2 du th.~1)
que l'intersection des noyaux des homomorphismes
$B'_\alpha\to B'_{\alpha\beta}$ (pour un $\alpha$ donné) est réduite à $0$~;
on en conclut que l'intersection des noyaux des homomorphismes composés
$v_{\alpha\beta}:B\to B'_\alpha\to B'_{\alpha\beta}$ ($\alpha$ et $\beta$
variables) est réduite à $0$. D'autre part, comme $B$ est entier sur $A$,
$B'_\alpha$ est entier sur $A'_\alpha$, donc ses idéaux maximaux sont au-dessus
de l'idéal maximal de $A'_\alpha$. Si l'on pose
$Z'_{\alpha\beta}=\operatorname{Spec}(B'_{\alpha\beta})$,
$T'_{\alpha\beta}=X'_\alpha\times_{Y'_\alpha}Z'_{\alpha\beta}$,
$f'_{\alpha\beta}=f_{(Z'_{\alpha\beta})}$,
$\mathcal G=\mathcal F\otimes_A B$ et
$\mathcal G'_{\alpha\beta}=\mathcal G\otimes_B B'_{\alpha\beta}$, on voit donc
que les hypothèses (i) et (ii) sont encore vérifiées lorsqu'on remplace
respectivement $A$, $A'_\alpha$, $u_\alpha$, $\mathcal F$ et $x$ par $B$,
$B'_{\alpha\beta}$, $v_{\alpha\beta}$, $\mathcal G$ et un point $t$ de
$T=X\otimes_A B$ au-dessus de $x$. Comme le corps résiduel de $B$ est
séparablement clos, on déduit de (11.4.11) que $\mathcal G$ est plat sur $B$ au
point $t$. Mais puisque $B$ est un $A$-module fidèlement plat, on conclut de
(2.1.4) que $\mathcal F$ est plat sur $A$ au point $x$, ce qui prouve (11.4.12).
\end{proof}

\begin{corollary}[11.4.13]
\label{IV.11.4.13-fr}
Soient $A$ un anneau artinien local, $Y=\operatorname{Spec}(A)$,
$f:X\to Y$ un morphisme localement de type fini, $\mathcal F$ un
$\mathcal O_X$-Module cohérent. Soit $(Y'_\alpha)$ une famille de
$Y$-préschémas, et pour tout $\alpha$, posons
$X'_\alpha=X\times_Y Y'_\alpha$, $f'_\alpha=f_{(Y'_\alpha)}$,
$\mathcal F'_\alpha=\mathcal F\otimes_Y Y'_\alpha$. Soit $x$ un point de $X$ et
supposons vérifiées les hypothèses suivantes~:
\begin{enumerate}[label=(\roman*)]
  \item L'intersection des noyaux des homomorphismes
  $A\to\Gamma(Y'_\alpha,\mathcal O_{Y'_\alpha})$ correspondant aux morphismes
  structuraux est réduite à $0$.
  \item Pour tout $\alpha$, $\mathcal F'_\alpha$ est $f'_\alpha$-plat en tous
  les points $x'_\alpha\in X'_\alpha$ dont la projection dans $X$ est $x$.
\end{enumerate}
Alors $\mathcal F$ est $f$-plat au point $x$.
\end{corollary}

\begin{proof}
En effet, pour tout $y'_\alpha\in Y'_\alpha$, considérons le schéma local
$Y''_{\alpha,y'_\alpha}=\operatorname{Spec}(\mathcal O_{Y'_\alpha,y'_\alpha})$~;
en vertu de (2.1.4),
$\mathcal F'_\alpha\otimes_{Y'_\alpha}Y''_{\alpha,y'_\alpha}$ est
$Y''_{\alpha,y'_\alpha}$-plat aux points dont les projections sur $X$ et
$Y''_{\alpha,y'_\alpha}$ sont $x$ et le point fermé de
$Y''_{\alpha,y'_\alpha}$. En outre le noyau de l'homomorphisme
$A\to\Gamma(Y'_\alpha,\mathcal O_{Y'_\alpha})$ est l'intersection des noyaux
des homomorphismes
$A\to\Gamma(Y''_{\alpha,y'_\alpha},\mathcal O_{Y''_{\alpha,y'_\alpha}})$, car
on se ramène aussitôt au cas où $Y'_\alpha$ est affine, et il suffit alors
d'appliquer Bourbaki, \emph{Alg. comm.}, chap.~II, \S~3, n$^\circ$~3, cor.~2 du
th.~1. En remplaçant la famille $(Y'_\alpha)$ par la famille des
$Y''_{\alpha,y'_\alpha}$, on est donc ramené à (11.4.12).
\end{proof}

\subsection{Descente de la platitude par des morphismes quelconques : cas général}
\label{subsection:IV.11.5-fr}

\begin{theorem}[11.5.1]
\label{IV.11.5.1-fr}
Soient $Y$ un préschéma localement noethérien, $f:X\to Y$ un morphisme
localement de type fini, $\mathcal F$ un $\mathcal O_X$-Module cohérent, $x$ un
point de $X$, $y=f(x)$. Soit $(Y'_\alpha)$ une famille de $Y$-préschémas locaux
$Y'_\alpha=\operatorname{Spec}(A'_\alpha)$ telle que les images des points
fermés $y'_\alpha$ des $Y'_\alpha$ soient toutes égales à $y$. Pour tout
$\alpha$, soient $\mathfrak m'_\alpha$ l'idéal maximal de $A'_\alpha$, et
$u_\alpha:\mathcal O_y\to A'_\alpha$
